\chapter{Tensores como Aplicaciones Multilineales}

La linealidad expresa que una aplicación respeta la estructura algebraica de
un espacio vectorial.
Sin embargo, muchas construcciones en matemáticas dependen de funciones que
toman varios argumentos vectoriales simultáneamente y permanecen lineales en
cada uno de ellos cuando los demás se mantienen fijos. Esto conduce de manera
natural a la noción de multilinealidad.

\medskip

Sean \(E_1, E_2, \dots, E_n\) y \(F\) espacios vectoriales sobre el mismo
cuerpo \(K\).

\begin{definition}[Aplicación multilineal]
Una aplicación
\[
T : E_1 \times E_2 \times \cdots \times E_n \to F
\]
se llama \emph{multilineal} si es lineal en cada argumento por separado.
Es decir, para cada índice \(i \in \{1,\dots,n\}\) y para todo
\[
x_i, y_i \in E_i, \quad \alpha \in K,
\]
se tiene
\[
T(x_1,\dots,x_i +_{E_i} y_i,\dots,x_n)
=
T(x_1,\dots,x_i,\dots,x_n)
+_F
T(x_1,\dots,y_i,\dots,x_n),
\]
y
\[
T(x_1,\dots,\alpha \circ x_i,\dots,x_n)
=
\alpha \circ T(x_1,\dots,x_i,\dots,x_n),
\]
con todos los demás argumentos fijos.
\end{definition}

Cuando \(n=1\), una aplicación multilineal es simplemente una aplicación lineal.
Así, las aplicaciones multilineales no son un tipo de objeto nuevo, sino una
generalización directa de las transformaciones lineales.

\medskip

Las aplicaciones multilineales aparecen de manera natural en muchos contextos:
productos bilineales, emparejamientos entre vectores y formas lineales,
e interacciones de orden superior que no pueden reducirse a una única entrada.
Su estudio sistemático conduce al concepto de tensor.

\medskip

Informalmente, un tensor es un objeto diseñado para codificar comportamiento
multilineal de manera intrínseca y sin coordenadas.

\medskip

A lo largo de este documento se utilizarán dos puntos de vista equivalentes.
El primero trata los tensores directamente como aplicaciones multilineales.

\medskip

Sea \(E\) un espacio vectorial sobre \(K\), y sea \(E^*\) su espacio dual,
es decir, el espacio vectorial de las aplicaciones lineales de \(E\) en \(K\).

\begin{definition}[Tensor como aplicación multilineal]
Un \emph{tensor de tipo \((k,\ell)\)} sobre \(E\) es una aplicación multilineal
\[
T : \underbrace{E^* \times \cdots \times E^*}_{k\ \text{veces}}
\times
\underbrace{E \times \cdots \times E}_{\ell\ \text{veces}}
\;\longrightarrow\; K.
\]
\end{definition}

Tal tensor toma \(k\) formas lineales y \(\ell\) vectores como argumentos,
y produce un escalar.
Se requiere linealidad en cada argumento por separado.

\medskip

Esta definición precisa varios casos familiares:
\begin{itemize}
\item Un escalar en \(K\) puede verse como un tensor de tipo \((0,0)\).
\item Un vector en \(E\) corresponde a un tensor de tipo \((0,1)\).
\item Una forma lineal en \(E^*\) es un tensor de tipo \((1,0)\).
\item Una forma bilineal sobre \(E\) es un tensor de tipo \((0,2)\).
\item Si $E$ es de dimensión finita, una aplicación lineal $T : E \to E$ puede
      identificarse con un tensor de tipo $(1,1)$.
\end{itemize}

Así, los vectores y las matrices reaparecen como los casos de orden más bajo
de un único marco unificado.

\medskip

A este nivel de abstracción no intervienen coordenadas, tablas ni
interpretaciones geométricas.
Un tensor queda definido enteramente por cómo actúa sobre vectores y formas
lineales, y por la multilinealidad de dicha acción.

% CORRECCIÓN 1: La motivación del producto tensorial está ahora claramente
% separada de su definición. La versión anterior mezclaba motivación y solución
% en el mismo bloque, con la frase final anticipando la definición que seguía
% a continuación (circular). La motivación es ahora autocontenida: identifica la
% limitación de las aplicaciones multilineales (no existe composición natural) y
% plantea la pregunta que el producto tensorial responderá, sin nombrar la
% solución prematuramente.

\medskip

La definición anterior es conceptualmente clara, pero tiene una limitación
práctica.
Las aplicaciones lineales se componen: si $S : E \to F$ y $T : F \to G$ son
lineales, entonces $T \circ S : E \to G$ es de nuevo lineal, y su matriz es
el producto de las matrices de $T$ y $S$.
Las aplicaciones multilineales no admiten una operación análoga.
Sus salidas y entradas no encajan de manera natural como para permitir una
composición canónica, y no existe ninguna regla general para combinar dos
aplicaciones multilineales en una tercera preservando la multilinealidad.

\medskip

Las aplicaciones multilineales forman espacios vectoriales bajo las operaciones
punto a punto, de modo que pueden sumarse y escalarse.
Lo que falta es una manera de \emph{reducir} la multilinealidad a la linealidad:
sustituir un problema multilineal por uno lineal sobre un espacio
convenientemente construido.
Esto es precisamente lo que logra el producto tensorial.

\medskip

Sean \(E_1,\dots,E_n\) espacios vectoriales sobre el mismo cuerpo \(K\).
Consideremos el conjunto de todas las aplicaciones multilineales
\[
T : E_1 \times \cdots \times E_n \longrightarrow F
\]
con valores en un espacio vectorial arbitrario \(F\).
El producto tensorial se define de modo que cada tal aplicación multilineal
corresponda de manera única a una aplicación lineal definida sobre un único
espacio vectorial.

\medskip

\begin{definition}[Producto tensorial]
El \emph{producto tensorial} de los espacios vectoriales \(E_1,\dots,E_n\) es
un espacio vectorial, denotado
\[
E_1 \otimes \cdots \otimes E_n,
\]
junto con una aplicación multilineal
\[
\otimes : E_1 \times \cdots \times E_n \longrightarrow
E_1 \otimes \cdots \otimes E_n,
\]
que satisface la siguiente propiedad universal:

Para todo espacio vectorial \(F\) y toda aplicación multilineal
\[
T : E_1 \times \cdots \times E_n \longrightarrow F,
\]
existe una única aplicación lineal
\[
\widetilde{T} : E_1 \otimes \cdots \otimes E_n \longrightarrow F
\]
tal que
\[
T(x_1,\dots,x_n) = \widetilde{T}(x_1 \otimes \cdots \otimes x_n)
\]
para todo \(x_i \in E_i\).
\end{definition}

\medskip

Esta propiedad caracteriza completamente el producto tensorial.
Los elementos \(x_1 \otimes \cdots \otimes x_n\) se llaman \emph{tensores simples}.
Todo elemento del espacio producto tensorial es una combinación lineal finita
de tensores simples.

\medskip

Como consecuencia, el estudio de las aplicaciones multilineales es equivalente
al estudio de las aplicaciones lineales definidas sobre productos tensoriales.
Concretamente,
\[
\text{Aplicaciones multilineales } E_1 \times \cdots \times E_n \to F
\quad\longleftrightarrow\quad
\text{Aplicaciones lineales } E_1 \otimes \cdots \otimes E_n \to F.
\]

\medskip

Esta equivalencia es el fundamento conceptual de la teoría de tensores.
Explica por qué los tensores pueden manipularse algebraicamente, componerse
con aplicaciones lineales y representarse concretamente una vez elegidas las bases.

\medskip

Volviendo al caso de un único espacio vectorial \(E\), un tensor de tipo
\((k,\ell)\) puede verse equivalentemente como una aplicación lineal
\[
T :
\underbrace{E^* \otimes \cdots \otimes E^*}_{k}
\otimes
\underbrace{E \otimes \cdots \otimes E}_{\ell}
\;\longrightarrow\; K.
\]

\medskip

Así los tensores son aplicaciones lineales definidas sobre productos tensoriales,
y su aparente complejidad surge únicamente al introducir coordenadas.

\medskip

A diferencia de la suma directa (Capítulo~I), que combina componentes
independientes de manera aditiva, el producto tensorial codifica interacciones
entre componentes mediante la multilinealidad.

\bigskip
\hrule
\medskip

\textit{Nota:} Una vez elegidas las bases, los productos tensoriales se
convierten en espacios vectoriales de dimensión finita, y los tensores admiten
representaciones coordenadas como tablas multidimensionales. Estas
representaciones dependen de las bases elegidas y no definen el tensor en sí.

\medskip
\hrule
\bigskip

Sea $E$ un espacio vectorial de dimensión finita sobre $K$, y sea
\[
B = \{ e_1, \dots, e_n \}
\]
una base de $E$. Su base dual
\[
B^* = \{ e^1, \dots, e^n \}
\subset E^*
\]
se define por las relaciones
\[
e^i(e_j) = \delta^i_j,
\]
donde $\delta^i_j$ es la delta de Kronecker.

\medskip

Todo vector $x \in E$ admite una expansión única
\[
x = x^i \circ e_i,
\]
y toda forma lineal $\varphi \in E^*$ admite una expansión única
\[
\varphi = \varphi_i \circ e^i.
\]

Los escalares $x^i$ y $\varphi_i$ son las \emph{componentes} de $x$ y de
$\varphi$ respecto a las bases elegidas.

\medskip

Sea ahora $T$ un tensor de tipo $(k,\ell)$ sobre $E$, es decir,
\[
T : \underbrace{E^* \times \cdots \times E^*}_{k}
\times
\underbrace{E \times \cdots \times E}_{\ell}
\longrightarrow K.
\]

La acción de $T$ queda completamente determinada por sus valores sobre los
elementos de la base:
\[
T(e^{i_1}, \dots, e^{i_k}, e_{j_1}, \dots, e_{j_\ell})
= T^{i_1 \dots i_k}{}_{j_1 \dots j_\ell}.
\]

Estos escalares se llaman las \emph{componentes de $T$ respecto a la base $B$}.
Codifican el tensor por completo una vez fijada una base.

\medskip

Para argumentos arbitrarios
\[
(\varphi^1, \dots, \varphi^k, x_1, \dots, x_\ell),
\]
la multilinealidad implica la expansión
\[
T(\varphi^1, \dots, \varphi^k, x_1, \dots, x_\ell)
=
T^{i_1 \dots i_k}{}_{j_1 \dots j_\ell}
\,
\varphi^1_{i_1} \cdots \varphi^k_{i_k}
\,
x_1^{j_1} \cdots x_\ell^{j_\ell}.
\]

Así, la notación de índices es la expresión coordenada de la multilinealidad
una vez elegida una base.

\medskip

Las aplicaciones bilineales proporcionan el ejemplo computacional más simple
de multilinealidad.
Una aplicación bilineal
\[
B : V \times W \to U
\]
es lineal en cada argumento por separado. Una vez elegidas bases en los espacios
vectoriales involucrados, la aplicación queda representada por coeficientes
$t_{ij}^{\;k}$ tales que las componentes del vector de salida satisfacen
\begin{equation}
z^{k} = \sum_{i=1}^{m}\sum_{j=1}^{n} t_{ij}^{\;k}\, x^{i} y^{j}.
\end{equation}

La colección de coeficientes $(t_{ij}^{\;k})$ forma un tensor de rango tres.
Esta representación muestra que los procedimientos computacionales familiares
son expresiones coordenadas de aplicaciones multilineales. Por ejemplo, la
multiplicación de números complejos
\begin{equation}
(x_1 + i x_2)(y_1 + i y_2)
= (x_1 y_1 - x_2 y_2)
+ i(x_1 y_2 + x_2 y_1)
\end{equation}
puede interpretarse como la evaluación de dos formas bilineales determinadas
por un tensor de coeficientes $2 \times 2 \times 2$.

Desde este punto de vista, los algoritmos que manipulan tablas numéricas son
implementaciones de transformaciones tensoriales escritas en un sistema de
referencia elegido. El cálculo hace por tanto explícita la estructura
multilineal invariante subyacente a las operaciones algebraicas.

\medskip

En particular, una aplicación bilineal
\[
B : E \times F \to G
\]
entre espacios vectoriales de dimensión finita queda representada, una vez
elegidas las bases, por componentes
\[
B^k{}_{ij},
\]
de modo que para vectores $x \in E$ e $y \in F$,
\[
(B(x,y))^k = B^k{}_{ij}\, x^i y^j.
\]
Así, las tablas multidimensionales de coeficientes que aparecen en los
algoritmos numéricos son representaciones coordenadas de tensores que codifican
operaciones multilineales.

\medskip

Una vez fijada una base, las componentes
\[
T^{i_1 \dots i_k}{}_{j_1 \dots j_\ell}
\]
pueden disponerse en una tabla multidimensional de escalares.

\medskip

Cambiar la base de $E$ cambia los valores numéricos de las componentes según
reglas de transformación precisas, mientras que el tensor en sí permanece
inalterado.

\bigskip
\hrule
\medskip

\textit{Nota: Tensores mediante leyes de transformación}

La definición de tensor como aplicación multilineal es intrínseca y libre de
coordenadas. Una caracterización equivalente pone de relieve cómo se
transforman las componentes de un tensor bajo un cambio de base.

\medskip

Este punto de vista hace explícito el contenido invariante que sobrevive a los
cambios de coordenadas, que es precisamente lo que permite que las leyes
tensoriales aparezcan en física y en aplicaciones de ingeniería con
independencia de la representación elegida.

\medskip

Sean $B = \{e_i\}$ y $B' = \{e'_i\}$ dos bases de $E$, relacionadas por
\[
e'_i = A^j{}_i \, e_j,
\]
donde $A^j{}_i$ es una matriz de cambio de base invertible.
Las bases duales correspondientes satisfacen
\[
e'^i = (A^{-1})^i{}_j \, e^j.
\]

\medskip

Sea $T$ un tensor de tipo $(k,\ell)$ sobre $E$, con componentes respecto a
$B$ dadas por
\[
T^{i_1 \dots i_k}{}_{j_1 \dots j_\ell}.
\]

La matriz $A^j{}_i$ expresa los nuevos vectores de la base en términos de los
anteriores, tal como se trató en el Capítulo~I.
En consecuencia, la regla de transformación siguiente describe cómo cambia la
representación coordenada de $T$ al pasar de la base $B$ a $B'$.

\medskip

Las componentes respecto a la nueva base $B'$ vienen dadas por
\[
T'^{i_1 \dots i_k}{}_{j_1 \dots j_\ell}
=
A^{i_1}{}_{p_1}
\cdots
A^{i_k}{}_{p_k}
\,
(A^{-1})^{q_1}{}_{j_1}
\cdots
(A^{-1})^{q_\ell}{}_{j_\ell}
\,
T^{p_1 \dots p_k}{}_{q_1 \dots q_\ell}.
\]

\medskip

Los índices contravariantes (índices superiores) se transforman con la matriz
$A$, mientras que los índices covariantes (índices inferiores) se transforman
con su inversa.
Esta regla de transformación no es una hipótesis adicional.
Es una consecuencia directa de la multilinealidad y de cómo se transforman
los vectores y los vectores duales bajo un cambio de base.

\medskip

Recíprocamente, supongamos que para cada elección de base se asigna una tabla
multidimensional de escalares
\[
T^{i_1 \dots i_k}{}_{j_1 \dots j_\ell}
\]
de tal manera que estas tablas están relacionadas por la ley de transformación
anterior siempre que se cambie la base.
Entonces estas colecciones coordenadas definen una única aplicación multilineal
\[
T :
\underbrace{E^* \times \cdots \times E^*}_{k}
\times
\underbrace{E \times \cdots \times E}_{\ell}
\longrightarrow K,
\]
independiente de la base elegida.

\medskip

Así los dos puntos de vista siguientes son equivalentes:
\begin{itemize}
\item Un tensor es una aplicación multilineal.
\item Un tensor es una colección de componentes que se transforman según la
      ley de transformación tensorial.
\end{itemize}

La primera definición enfatiza la estructura intrínseca.
La segunda enfatiza la invariancia bajo cambio de sistema de referencia.
Su equivalencia explica por qué la teoría de tensores desempeña un papel
central en geometría y física: las ecuaciones tensoriales son válidas en todo
sistema de coordenadas, porque ambos lados se transforman según la misma ley.

\medskip
\hrule
\bigskip

Así:
\begin{itemize}
\item los tensores son objetos multilineales abstractos;
\item las tablas son representaciones dependientes de la base;
\item los índices registran cómo la multilinealidad se descompone en coordenadas.
\end{itemize}

Esta distinción es esencial para interpretar las implementaciones computacionales.

\medskip

Hasta este punto, los tensores se han considerado como objetos multilineales
estáticos.
Introducimos ahora la operación que les permite interactuar algebraicamente.

% CORRECCIÓN 3: La frase rota («operación. introducimos ahora...») y la
% repetición de la observación sobre la no composición han sido eliminadas.
% La transición es ahora una única frase limpia que abre la sección de contracción.

\medskip

Para combinar tensores y producir nuevos objetos multilineales, se introduce
la operación fundamental de contracción tensorial, que generaliza el producto
de matrices y explica la aparición de sumas sobre índices repetidos.

\medskip

Denotamos por
\[
\mathcal{T}^{(k,\ell)}(E)
\]
el espacio vectorial de todos los tensores de tipo $(k,\ell)$ sobre $E$.

\medskip

Sea $E$ un espacio vectorial de dimensión finita sobre $K$.
Consideremos un tensor
\[
T \in \mathcal{T}^{(k,\ell)}(E)
\]
y un tensor
\[
S \in \mathcal{T}^{(m,n)}(E).
\]

Supongamos que se selecciona un índice covariante de $T$ y un índice
contravariante de $S$.
La contracción tensorial es la operación que empareja estos dos índices mediante
la aplicación de evaluación natural
\[
E^* \times E \to K.
\]

\begin{definition}[Contracción tensorial]
Sean
\[
T \in \mathcal{T}^{(k,\ell)}(E),
\qquad
S \in \mathcal{T}^{(m,n)}(E).
\]
Eligiendo un índice covariante de $T$ y un índice contravariante de $S$, su
contracción se define componiendo estos argumentos con el emparejamiento de
evaluación natural
\[
\langle\cdot,\cdot\rangle : E^* \times E \to K.
\]
En coordenadas, esta operación aparece como una suma sobre una base.

\medskip

El tensor resultante tiene tipo
\[
(k+m-1,\;\ell+n-1).
\]
\end{definition}

\medskip

La contracción reduce el número total de índices del tensor en dos: un índice
contravariante y uno covariante se emparejan y eliminan.
No se requieren coordenadas para definir esta operación.

\medskip

Para expresar la contracción en coordenadas, debe elegirse una base.
El emparejamiento abstracto aparece entonces como una suma sobre índices repetidos.

\medskip

Fijemos una base $\{e_i\}$ de $E$ y su base dual $\{e^i\}$.

Sean
\[
T^{i}{}_{j}
\quad\text{y}\quad
S^{j}{}_{k}
\]
las componentes de dos tensores de tipo $(1,1)$.

\medskip

La contracción sobre el índice $j$ produce un nuevo tensor con componentes
\[
R^{i}{}_{k}
=
T^{i}{}_{j} \, S^{j}{}_{k}.
\]

El índice repetido $j$ indica una suma:
\[
R^{i}{}_{k}
=
\sum_{j=1}^{n}
T^{i}{}_{j} \, S^{j}{}_{k}.
\]
Esta suma no es una regla adicional.
Es la expresión coordenada del emparejamiento abstracto entre $E^*$ y $E$.

\medskip

Sean $A : E \to E$ y $B : E \to E$ aplicaciones lineales.
Cada una puede identificarse con un tensor de tipo $(1,1)$.
Respecto a una base, sus componentes se escriben como
\[
A^{i}{}_{j},
\qquad
B^{j}{}_{k}.
\]

La composición $A \circ B : E \to E$ tiene componentes
\[
(A \circ B)^{i}{}_{k}
=
A^{i}{}_{j} \, B^{j}{}_{k}.
\]

Así, el producto de matrices es exactamente la contracción tensorial entre dos
tensores de tipo $(1,1)$, y la regla de suma familiar para el producto de
matrices es un caso especial del emparejamiento abstracto entre $E^*$ y $E$.

\medskip

A nivel abstracto tenemos:
\[
\text{Contracción tensorial}
\]
que, tras elegir bases, se convierte en
\[
\text{Suma sobre índices}
\quad\longrightarrow\quad
C^{i}{}_{k} = \sum_j A^{i}{}_{j} B^{j}{}_{k}.
\]

% CORRECCIÓN 4: El párrafo que comenzaba con «Es fundamental reconocer...»
% ha sido eliminado. Su tono era marcadamente más retórico que el texto
% circundante, introducía el término no estándar «demolición de índices»,
% y su contenido (la contracción como motor del cálculo) ya se transmite
% de manera más precisa en la nota sobre Arrays Sistólicos más adelante.

% CORRECCIÓN 2: El ejemplo de la cuadrícula discreta (tabla de fuerzas) se ha
% trasladado aquí, tras definir la contracción. El lector dispone ahora de las
% herramientas algebraicas para comprender qué significan los índices de la tabla
% y cómo podrían contraerse. Anteriormente aparecía antes de la contracción,
% dejando la tabla sin ningún contexto operacional.

\medskip

En contextos computacionales y de modelado, los tensores aparecen con frecuencia
a través de tablas multidimensionales cuyos índices etiquetan parámetros
independientes del problema.
El siguiente ejemplo ilustra cómo surgen en la práctica tales tablas mientras
siguen representando datos tensoriales.

\medskip

Consideremos un sistema físico modelado sobre una cuadrícula bidimensional
discreta.
En cada punto espacial $(i,j)$ y en cada instante de tiempo $t$, se mide un
vector de fuerza
\[
F(t,i,j) \in \mathbb{R}^2
\]
con componentes $(F_x, F_y)$ respecto a un sistema de referencia espacial fijo.

\medskip

Sean:
\begin{itemize}
\item $t \in \{1,\dots,T\}$ el índice de tiempo,
\item $(i,j) \in \{1,\dots,N\} \times \{1,\dots,M\}$ el índice de posición espacial,
\item $\alpha \in \{1,2\}$ el índice de componentes de la fuerza.
\end{itemize}

Los datos recopilados pueden escribirse como una tabla
\[
F_{t i j \alpha},
\]
de forma $(T,N,M,2)$.

\medskip

Cada índice corresponde a una elección de modelado independiente: tiempo,
espacio y componentes vectoriales.
La tabla almacena valores numéricos respecto a sistemas de coordenadas elegidos.
Solo el índice de componente $\alpha$ corresponde a un índice vectorial genuino
en el sentido tensorial; los índices $(t,i,j)$ etiquetan parámetros
independientes del modelo.
Contraer el índice $\alpha$ con una forma lineal produce, en cada punto del
espacio-tiempo, una medición escalar de la fuerza en una dirección elegida.

\medskip

Abstractamente, estos datos pueden interpretarse como la representación
coordenada de una función con valores tensoriales sobre un dominio discreto.

\bigskip
\hrule
\medskip

\textit{Nota: El álgebra lineal realizada en hardware --- Arrays sistólicos
\cite{kung}}

\medskip

La interpretación del producto de matrices como contracción tensorial no es
solo una reformulación conceptual.
También determina cómo se diseña el hardware computacional moderno.

\medskip

En los grandes cálculos numéricos, la operación dominante es la evaluación
repetida de productos internos,
\[
c_i = \sum_{j=1}^{n} a_{ij} b_j,
\]
que es la expresión coordenada de una aplicación lineal aplicada a un vector.
Equivalentemente, es una contracción entre un tensor de tipo $(1,1)$ y un
tensor de tipo $(0,1)$.

\medskip

A finales de la década de 1970, H.~T.~Kung y C.~E.~Leiserson propusieron
organizar el cálculo en torno a esta operación mediante arquitecturas llamadas
\emph{arrays sistólicos}.
En lugar de transferir datos de ida y vuelta entre la memoria y un procesador
central, los datos fluyen rítmicamente a través de una cuadrícula de elementos
de procesamiento simples.
Cada elemento realiza repetidamente una única operación de la forma
\[
C \leftarrow A B + C,
\]
un paso de multiplicación-acumulación que implementa una contribución a un
producto interno.

\medskip

Matemáticamente, cada elemento de procesamiento evalúa una contracción tensorial
parcial.
A medida que los datos se propagan por el array, las contracciones sucesivas
se acumulan, y el resultado global emerge de la acción coordinada de muchas
operaciones lineales locales.

\medskip

Esta organización refleja un principio estructural: la eficiencia del hardware
se logra cuando el cálculo refleja la estructura algebraica.
Dado que las transformaciones lineales se descomponen en sumas de productos,
una arquitectura optimizada para las operaciones de multiplicación-acumulación
implementa de manera natural el álgebra lineal a gran escala.

\medskip

Las arquitecturas computacionales modernas de alto rendimiento adoptan con
frecuencia variantes de este paradigma, disponiendo muchos elementos de
procesamiento simples para evaluar operaciones de multiplicación-acumulación
en paralelo.
Lo que aparece a nivel de software como contracción tensorial o producto de
matrices se ejecuta físicamente como flujos de productos internos coordinados.

\medskip

Desde el presente punto de vista, este desarrollo no es accidental.
La contracción tensorial es la interacción fundamental entre objetos
multilineales, y el cálculo a gran escala reduce los algoritmos complejos a
evaluaciones repetidas de esta interacción.

\medskip

En este sentido, las tablas numéricas almacenadas en memoria, las sumas sobre
índices escritas en fórmulas y los datos que fluyen por circuitos de silicio
son tres realizaciones del mismo objeto matemático: la evaluación de
aplicaciones multilineales.

\medskip
\hrule
\bigskip

Los principios abstractos desarrollados anteriormente aparecen de manera
concreta en la física clásica, donde las cantidades tensoriales describen
relaciones invariantes entre observables físicas.

\medskip

\textit{Un ejemplo clásico: el tensor de inercia \cite{feynman}}

En mecánica clásica, la distribución de masa en un cuerpo rígido se codifica
mediante un tensor de segundo orden llamado \emph{tensor de inercia}.
Respecto a una base ortonormal elegida de $\mathbb{R}^3$, sus componentes
vienen dadas por
\[
I_{ij} = \int_{\Omega}
\big( \delta_{ij} \|x\|^2 - x_i x_j \big)\, dm,
\]
donde $\Omega$ es el cuerpo y $x$ es el vector de posición.

\medskip

El tensor de inercia define una relación lineal entre la velocidad angular
$\omega$ y el momento angular $L$:
\[
L_i = I_{ij} \, \omega^j.
\]

\medskip

Bajo un cambio de base determinado por una matriz invertible $A$, sus
componentes se transforman según
\[
I'_{ij}
=
A^{p}{}_{i} \,
A^{q}{}_{j} \,
I_{pq},
\]
que es precisamente la ley de transformación de un tensor de tipo $(0,2)$.

\medskip

Así, aunque las entradas numéricas de $I_{ij}$ dependen del sistema de
referencia elegido, el objeto geométrico que representa no depende de él.
El contenido físico ---la resistencia a la aceleración rotacional--- es
invariante bajo cambio de coordenadas.

\medskip

La teoría de tensores extiende por tanto el principio central introducido en
el Capítulo~I: la estructura matemática se identifica no por su representación
numérica, sino por la invariancia de sus relaciones definidoras bajo las
transformaciones admisibles.

% CORRECCIÓN 5 y 6: Las notas se reordenan para agrupar los ejemplos aplicados
% y separarlos de la perspectiva especulativa.
% Orden: Perspectivas (hardware futuro) -> observación sobre chi-cuadrado
% (puente con el Cap. II) -> cubos OLAP (aplicación en análisis de datos)
% -> Arrays sistólicos (hardware).
% El párrafo introductorio redundante que precedía a la nota sobre Arrays
% Sistólicos en el texto principal ha sido eliminado; la apertura de la propia
% nota resulta ahora suficiente.

\section*{Notas y Observaciones}

\bigskip
\hrule
\medskip

\textit{Perspectiva: Computación más allá de la Electrónica \cite{Xu}}

\medskip

Los avances recientes en arquitecturas de cómputo exploran sustratos físicos
más allá de los circuitos electrónicos tradicionales, entre ellos los sistemas
ópticos y fotónicos en los que la información se procesa mediante
interferencia, difracción y propagación de la luz.

\medskip

A pesar de sus diferencias físicas, estos sistemas están concebidos para
implementar las mismas operaciones algebraicas fundamentales que aparecen a lo
largo de este capítulo.
En particular, la tarea computacional dominante sigue siendo la evaluación de
expresiones de la forma
\[
C^{i}{}_{k} = A^{i}{}_{j} B^{j}{}_{k},
\]
es decir, la contracción tensorial.

\medskip

En estas arquitecturas emergentes, la interacción de las señales físicas
realiza de forma natural operaciones de multiplicación y acumulación en
paralelo, de modo que la evaluación de aplicaciones multilineales es ejecutada
directamente por el proceso físico subyacente, sin necesidad de instrucciones
electrónicas secuenciales.

\medskip

Sin embargo, surge una restricción no trivial debida a la codificación física
de la información.
Los sistemas ópticos no coherentes representan valores como intensidades
luminosas, que son intrínsecamente no negativas; operan, por tanto, dentro del
dominio restringido~$\mathbb{R}^+$ en lugar de sobre la recta real completa.
Dado que las contracciones tensoriales exigen pesos y activaciones que recorren
todo~$\mathbb{R}$, tales sistemas resultan algebraicamente incompletos para el
cómputo en general.
Las arquitecturas coherentes resuelven esta limitación codificando los valores
como \emph{amplitudes} de campo óptico --- magnitudes que admiten signo --- y
recuperando la salida con signo mediante interferencia en la detección.
El grupo completo $(\mathbb{R},+)$ queda así disponible, y las contracciones
tensoriales de signo arbitrario son realizables directamente en el hardware.

\medskip

Desde esta perspectiva, los avances en hardware no alteran la estructura
matemática de la computación.
Proporcionan nuevas realizaciones físicas de las mismas operaciones
invariantes: transformaciones lineales y contracciones tensoriales, ahora
extendidas a su dominio natural.
La completitud algebraica que aquí se requiere --- el paso
de~$\mathbb{R}^+$ a~$\mathbb{R}$ --- es precisamente la distinción estructural
entre un grupo multiplicativo y uno aditivo, cuya correspondencia es el objeto
del Capítulo~IV.

\medskip

Así, con independencia de si el cómputo se lleva a cabo por medios
electrónicos, ópticos o de cualquier otra naturaleza física, el marco
algebraico desarrollado aquí permanece inalterado.

\medskip
\hrule
\bigskip

\bigskip
\hrule
\medskip

\textit{Observación: Puente con el Capítulo~II --- el estadístico chi-cuadrado
como tensor de tipo $(0,2)$}

\medskip

El estadístico chi-cuadrado encontrado en el Capítulo~II proporciona un ejemplo
anterior de tensor de tipo $(0,2)$.
Recordemos que la expresión
\[
V
=
\sum_{i=1}^{n}
\frac{(Y_i - np_i)^2}{np_i}
\]
fue identificada como la norma al cuadrado inducida por una forma bilineal
simétrica definida positiva $B$ sobre $\mathbb{R}^n$, cuya matriz en la base
estándar es diagonal con entradas $(np_i)^{-1}$.
Desde el presente punto de vista, $B$ es precisamente un tensor de tipo $(0,2)$:
una aplicación multilineal
\[
B : \mathbb{R}^n \times \mathbb{R}^n \longrightarrow \mathbb{R}
\]
que es simétrica y definida positiva.
Su representación coordenada cambia con la base, pero el objeto geométrico que
define ---los conjuntos de nivel elipsoidales de $V$--- no cambia.
Así, la cantidad estadística $V$, la forma cuadrática del Capítulo~II y el
tensor de tipo $(0,2)$ del presente capítulo son tres descripciones de la
misma estructura invariante.

\medskip
\hrule
\bigskip

\bigskip
\hrule
\medskip

\textit{Nota: Cubos OLAP y contracción tensorial}

\medskip

En el análisis de datos a gran escala, las observaciones se organizan con
frecuencia a lo largo de múltiples dimensiones categóricas independientes como
tiempo, ubicación, componente del sistema o tipo de medición.
La estructura resultante se llama comúnmente \emph{cubo OLAP} (cubo de
Procesamiento Analítico en Línea, por sus siglas en inglés).

\medskip

Matemáticamente, tal objeto asigna valores numéricos a elementos de un
producto cartesiano
\[
D_1 \times D_2 \times \cdots \times D_n,
\]
donde cada conjunto $D_i$ representa una dimensión de observación.
Tras enumerar estas dimensiones, los datos pueden escribirse como una tabla
multidimensional
\[
F_{i_1 i_2 \dots i_n},
\]
que es precisamente la representación coordenada de un tensor indexado por
$n$ parámetros independientes.

\medskip

En la práctica, solo se observa un pequeño subconjunto de todas las
combinaciones posibles de índices, por lo que la tabla es típicamente dispersa.
Esto refleja de nuevo una distinción enfatizada a lo largo de este capítulo:
las tablas dependen de las coordenadas y las convenciones de almacenamiento
elegidas, mientras que la estructura multilineal subyacente es independiente
de la representación.

\medskip

Las consultas analíticas suelen agregar datos a lo largo de dimensiones
seleccionadas.
Por ejemplo, combinar observaciones ignorando ciertos índices produce cantidades
de la forma
\[
G_{i_1 \dots i_k}
=
\sum_{j_1,\dots,j_m}
F_{i_1 \dots i_k j_1 \dots j_m}.
\]
Desde el punto de vista tensorial, esta operación es una contracción: los
índices correspondientes a las dimensiones agregadas se suman, reduciendo el
orden del tensor.

\medskip

Así, las operaciones conocidas operacionalmente como filtrado, agrupación o
agregación corresponden matemáticamente a proyecciones de un tensor de alta
dimensión sobre subtensores de dimensión inferior.
El cubo OLAP proporciona por tanto una realización concreta de la estructura
tensorial que surge en el análisis de datos moderno.

\medskip

Este ejemplo refuerza un tema recurrente: la contracción tensorial aparece
dondequiera que las relaciones multidimensionales se resumen mediante
agregación sistemática.

\medskip
\hrule
\bigskip